\documentclass[10pt,a4paper]{amsart}
\usepackage{amsmath,amssymb,amsthm,graphicx}
\usepackage{verbatim}
\newenvironment{reference}{\comment}{\endcomment}
%\newenvironment{reference}{}{}
\newenvironment{slogan}{\comment}{\endcomment}
\newenvironment{history}{\comment}{\endcomment}

% For commutative diagrams we use Xy-pic
\usepackage[all]{xy}

% We use 2cell for 2-commutative diagrams.
\xyoption{2cell}
% Theorem environments.
%
\theoremstyle{plain}
\newtheorem{theorem}[subsection]{Theorem}
\newtheorem{proposition}[subsection]{Proposition}
\newtheorem{lemma}[subsection]{Lemma}

\theoremstyle{definition}
\newtheorem{definition}[subsection]{Definition}
\newtheorem{example}[subsection]{Example}
\newtheorem{exercise}[subsection]{Exercise}
\newtheorem{situation}[subsection]{Situation}

\theoremstyle{remark}
\newtheorem{remark}[subsection]{Remark}
\newtheorem{remarks}[subsection]{Remarks}

\numberwithin{equation}{subsection}

% Macros
%
\def\lim{\mathop{\mathrm{lim}}\nolimits}
\def\colim{\mathop{\mathrm{colim}}\nolimits}
\def\Spec{\mathop{\mathrm{Spec}}}
\def\Hom{\mathop{\mathrm{Hom}}\nolimits}
\def\Ext{\mathop{\mathrm{Ext}}\nolimits}
\def\SheafHom{\mathop{\mathcal{H}\!\mathit{om}}\nolimits}
\def\SheafExt{\mathop{\mathcal{E}\!\mathit{xt}}\nolimits}
\def\Sch{\mathit{Sch}}
\def\Mor{\mathop{\mathrm{Mor}}\nolimits}
\def\Ob{\mathop{\mathrm{Ob}}\nolimits}
\def\Sh{\mathop{\mathit{Sh}}\nolimits}
\def\NL{\mathop{N\!L}\nolimits}
\def\CH{\mathop{\mathrm{CH}}\nolimits}
\def\proetale{{pro\text{-}\acute{e}tale}}
\def\etale{{\acute{e}tale}}
\def\QCoh{\mathit{QCoh}}
\def\Ker{\mathop{\mathrm{Ker}}}
\def\Im{\mathop{\mathrm{Im}}}
\def\Coker{\mathop{\mathrm{Coker}}}
\def\Coim{\mathop{\mathrm{Coim}}}

% Boxtimes
%
\DeclareMathSymbol{\boxtimes}{\mathbin}{AMSa}{"02}

%
% Macros for moduli stacks/spaces
%
\def\QCohstack{\mathcal{QC}\!\mathit{oh}}
\def\Cohstack{\mathcal{C}\!\mathit{oh}}
\def\Spacesstack{\mathcal{S}\!\mathit{paces}}
\def\Quotfunctor{\mathrm{Quot}}
\def\Hilbfunctor{\mathrm{Hilb}}
\def\Curvesstack{\mathcal{C}\!\mathit{urves}}
\def\Polarizedstack{\mathcal{P}\!\mathit{olarized}}
\def\Complexesstack{\mathcal{C}\!\mathit{omplexes}}
% \Pic is the operator that assigns to X its picard group, usage \Pic(X)
% \Picardstack_{X/B} denotes the Picard stack of X over B
% \Picardfunctor_{X/B} denotes the Picard functor of X over B
\def\Pic{\mathop{\mathrm{Pic}}\nolimits}
\def\Picardstack{\mathcal{P}\!\mathit{ic}}
\def\Picardfunctor{\mathrm{Pic}}
\def\Deformationcategory{\mathcal{D}\!\mathit{ef}}
\usepackage{cleveref}
\crefname{originallemma}{original lemma}{original lemmas}
\Crefname{originallemma}{Original lemma}{Original lemmas}
\makeatletter
\newlabel{algebra-lemma-exists-dvr}{{00PH}{source}{}{sourcepointer-1}{}}
\newlabel{algebra-lemma-exists-dvr@cref}{{[originallemma][0][]00PH}{[1][0][]source}{}{}{}}
\makeatother
\crefname{originallemma}{original lemma}{original lemmas}
\Crefname{originallemma}{Original lemma}{Original lemmas}
\makeatletter
\newlabel{formal-spaces-lemma-adic-into-completion}{{0GBU}{source}{}{sourcepointer-2}{}}
\newlabel{formal-spaces-lemma-adic-into-completion@cref}{{[originallemma][0][]0GBU}{[1][0][]source}{}{}{}}
\makeatother
\crefname{originallemma}{original lemma}{original lemmas}
\Crefname{originallemma}{Original lemma}{Original lemmas}
\makeatletter
\newlabel{formal-spaces-lemma-formal-completion-types}{{0AQ1}{source}{}{sourcepointer-3}{}}
\newlabel{formal-spaces-lemma-formal-completion-types@cref}{{[originallemma][0][]0AQ1}{[1][0][]source}{}{}{}}
\makeatother
\crefname{originalreference}{original reference}{original references}
\Crefname{originalreference}{Original reference}{Original references}
\makeatletter
\newlabel{lemma-map-completions-finite-type}{{0AQ9}{source}{}{sourcepointer-4}{}}
\newlabel{lemma-map-completions-finite-type@cref}{{[originalreference][0][]0AQ9}{[1][0][]source}{}{}{}}
\makeatother
\crefname{originallemma}{original lemma}{original lemmas}
\Crefname{originallemma}{Original lemma}{Original lemmas}
\makeatletter
\newlabel{more-algebra-lemma-completion-dvr}{{0AP1}{source}{}{sourcepointer-5}{}}
\newlabel{more-algebra-lemma-completion-dvr@cref}{{[originallemma][0][]0AP1}{[1][0][]source}{}{}{}}
\makeatother
\crefname{originallemma}{original lemma}{original lemmas}
\Crefname{originallemma}{Original lemma}{Original lemmas}
\makeatletter
\newlabel{spaces-morphisms-lemma-characterize-proper}{{0A40}{source}{}{sourcepointer-6}{}}
\newlabel{spaces-morphisms-lemma-characterize-proper@cref}{{[originallemma][0][]0A40}{[1][0][]source}{}{}{}}
\makeatother
\title{The completion morphism is rig-surjective when f is proper and U' to U surjective}
\author{The Stacks Project Authors}
\date{stacks-project:restricted.tex:local-snapshot-20261012; GNU FDL 1.2 or later}
\begin{document}
\maketitle
\noindent\textit{Editorial note.} All delivered mathematical text is the author's own, copied from the pinned Stacks Project chapter file restricted.tex; no statement or proof step was rewritten, reordered or removed.\par\medskip
\noindent\textit{Editorial note.} The counted proof is the author's complete argument: the proof algebraizes the adic test morphism, obtains a K'-point of U' over Spec(K) from surjectivity and properness, dominates R by a DVR with fraction field K', extends over the closed point by the valuative criterion of properness, and observes that T' pulls back to the closed point, finishing with the completed DVR lift\par\medskip
\noindent\textit{Editorial note.} The label lemma-completion-proper-surjective-rig-surjective names the counted statement, and the proof environment immediately following it in the source is delivered verbatim as its proof.\par\medskip
\noindent\textit{Editorial note.} The counted statement is carried by the author's own lemma environment in the chapter source.\par\medskip
\noindent\textit{Allowed standard prerequisites (editorial).}\begin{itemize}
\item the completion-functor situation of the section: f : X' -> X with T' the inverse image of T, U' = f\textasciicircum{}\{-1\}U, and the morphism of formal completions f\_\{/T\} : X'\_\{/T'\} -> X\_\{/T\} with its cartesian diagram
\item adic morphisms from Spf(R) of a complete discrete valuation ring into a formal completion and their algebraization to morphisms of schemes (Formal Spaces, Lemma formal-spaces-lemma-adic-into-completion)
\item the valuative criterion of properness for algebraic spaces (Morphisms of Spaces, Lemma spaces-morphisms-lemma-characterize-proper)
\item existence of a discrete valuation ring R' with fraction field K' dominating a given DVR R (Algebra, Lemma algebra-lemma-exists-dvr)
\item the completion of a discrete valuation ring is a complete discrete valuation ring (More on Algebra, Lemma more-algebra-lemma-completion-dvr)
\item rig-surjectivity of a morphism of formal algebraic spaces: every adic test morphism from a complete DVR lifts after passing to an extension
\end{itemize}\medskip
\ifdefined\hypertarget\hypertarget{sourcepointer-1}{}\fi
\noindent\textit{Source reference.} Original lemma 00PH: sources/stacks-chow/tags.txt, source lines 931--931. Prerequisite pointer invoked at proof line(s) 6227; it names a standard fact of a sibling chapter that the counted argument uses, and every step of the counted proof built on it is retained verbatim in the delivered body.\par\medskip
\ifdefined\hypertarget\hypertarget{sourcepointer-2}{}\fi
\noindent\textit{Source reference.} Original lemma 0GBU: sources/stacks-chow/tags.txt, source lines 20384--20384. Prerequisite pointer invoked at proof line(s) 6213; it names a standard fact of a sibling chapter that the counted argument uses, and every step of the counted proof built on it is retained verbatim in the delivered body.\par\medskip
\ifdefined\hypertarget\hypertarget{sourcepointer-3}{}\fi
\noindent\textit{Source reference.} Original lemma 0AQ1: sources/stacks-chow/tags.txt, source lines 13494--13494. Prerequisite pointer invoked at proof line(s) 6209; it names a standard fact of a sibling chapter that the counted argument uses, and every step of the counted proof built on it is retained verbatim in the delivered body.\par\medskip
\ifdefined\hypertarget\hypertarget{sourcepointer-4}{}\fi
\noindent\textit{Source reference.} Original reference 0AQ9: sources/stacks-restricted/pinned-original-source.tex, source lines 6009--6009. Prerequisite pointer invoked at proof line(s) 6208; it names a standard fact of the same chapter that the counted argument uses, and every step of the counted proof built on it is retained verbatim in the delivered body.\par\medskip
\ifdefined\hypertarget\hypertarget{sourcepointer-5}{}\fi
\noindent\textit{Source reference.} Original lemma 0AP1: sources/stacks-chow/tags.txt, source lines 13459--13459. Prerequisite pointer invoked at proof line(s) 6237; it names a standard fact of a sibling chapter that the counted argument uses, and every step of the counted proof built on it is retained verbatim in the delivered body.\par\medskip
\ifdefined\hypertarget\hypertarget{sourcepointer-6}{}\fi
\noindent\textit{Source reference.} Original lemma 0A40: sources/stacks-chow/tags.txt, source lines 12758--12758. Prerequisite pointer invoked at proof line(s) 6230; it names a standard fact of a sibling chapter that the counted argument uses, and every step of the counted proof built on it is retained verbatim in the delivered body.\par\medskip
\begin{lemma}
\label{lemma-completion-proper-surjective-rig-surjective}
In the situation above. If $X$ is locally Noetherian,
$f$ is proper, and $U' \to U$ is surjective, then $f_{/T}$ is rig-surjective.
\end{lemma}

\begin{proof}
(The statement makes sense by
Lemma \ref{lemma-map-completions-finite-type} and
Formal Spaces, Lemma \ref{formal-spaces-lemma-formal-completion-types}.)
Let $R$ be a complete discrete valuation ring with fraction field $K$.
Let $p : \text{Spf}(R) \to X_{/T}$ be an adic morphism of
formal algebraic spaces. By Formal Spaces, Lemma
\ref{formal-spaces-lemma-adic-into-completion}
the composition $\text{Spf}(R) \to X_{/T} \to X$
corresponds to a morphism $q : \Spec(R) \to X$
which maps $\Spec(K)$ into $U$. Since $U' \to U$ is proper and surjective
we see that $\Spec(K) \times_U U'$ is nonempty and proper over $K$.
Hence we can choose a field extension $K'/K$ and a commutative
diagram
$$
\xymatrix{
\Spec(K') \ar[r] \ar[d] & U' \ar[r] \ar[d] & X' \ar[d] \\
\Spec(K) \ar[r] & U \ar[r] & X
}
$$
Let $R' \subset K'$ be a discrete valuation ring dominating $R$
with fraction field $K'$, see Algebra, Lemma \ref{algebra-lemma-exists-dvr}.
Since $\Spec(K) \to X$ extends to $\Spec(R) \to X$ we see by the valuative
criterion of properness
(Morphisms of Spaces, Lemma \ref{spaces-morphisms-lemma-characterize-proper})
that we can extend our $K'$-point of $U'$ to a morphism
$\Spec(R') \to X'$ over $\Spec(R) \to X$.
It follows that the inverse image of $T'$ in $\Spec(R')$ is the
closed point and we find an adic morphism
$\text{Spf}((R')^\wedge) \to X'_{/T'}$ lifting $p$
as desired (note that $(R')^\wedge$ is a complete discrete valuation ring
by More on Algebra, Lemma \ref{more-algebra-lemma-completion-dvr}).
\end{proof}

\end{document}
